%============================================================
% Chapter 25: Minimum Cut & Flow Modelling
% STATUS: PLANNED  (generated by tools/outline.py; replace this file when the chapter is written)
%============================================================
\chapter{Minimum Cut \& Flow Modelling}
\label{ch:min_cut}

\begin{tcolorbox}[enhanced, colback=white, colframe=secondaryblue, boxrule=1pt, arc=2mm]
\textbf{\large\textcolor{primaryblue}{Every maximum flow certifies a minimum cut --- and cuts model selection, separation and closure problems.}}
\end{tcolorbox}

\chapterfacts{2000}{2700}{5}{10}{30}

\begin{plannedbox}[Planned Contents --- this chapter is not written yet]
\textbf{Part A --- Lessons}
\begin{enumerate}[leftmargin=*, label=\textbf{A\arabic*.}, itemsep=1pt]
    \item The max-flow min-cut theorem and recovering the cut
    \item Project selection and maximum-weight closure
    \item Infinite edges as constraints
    \item Flows with lower bounds; circulations
    \item Cuts in grids and images
\end{enumerate}
\textbf{Part B --- Patterns} (with candidate problems, checked when the chapter is written)
\begin{enumerate}[leftmargin=*, label=\textbf{P\arabic*.}, itemsep=1pt]
    \item \textbf{Recovering the cut} --- Police Chase {\small\color{darkgray}(CSES 1695)}
    \item \textbf{Project selection} --- Petya and Graph {\small\color{darkgray}(CF 1082G)}
    \item \textbf{Lower bounds and circulations}
    \item \textbf{Grid cuts}
    \item \textbf{Cut modelling}
\end{enumerate}
\end{plannedbox}
